% INFOGAIN -- patient-level information-value decomposition of multimodal testing
% Target: Nature Communications / npj Digital Medicine
% Build:  latexmk -pdf main.tex   (requires numbers.tex, generated by
%         scripts/make_paper_numbers.py from the results tree)
\documentclass[11pt]{article}

\usepackage[margin=1in]{geometry}
\usepackage{amsmath,amssymb,amsthm,mathtools}
\usepackage{graphicx}
\usepackage{booktabs}
\usepackage{authblk}
\usepackage[hidelinks]{hyperref}
\usepackage[numbers,sort&compress]{natbib}
\usepackage{xcolor}
\usepackage{caption}
\usepackage{siunitx}
% CM Type1 fonts in a minimal TeX Live are not scalable, so font expansion
% cannot run; protrusion (the part that actually matters here) still does.
\usepackage[expansion=false]{microtype}

% Cross-document references into the supplement. Hand-numbered "Table~S7"
% citations had already drifted twice from the actual ordering, silently; xr
% makes a stale reference render as ?? instead of pointing at the wrong table.
% Requires supplementary.aux, so build the supplement first (see paper/README).
\usepackage{xr}
\externaldocument{supplementary}

\captionsetup{font=small,labelfont=bf}
\graphicspath{{figures/}{../results/cohort_synthetic/mortality_30d/figures/}%
              {../results/simulation/figures/}}

% Wide generated tables: shrink to the text block only when they overrun it.
% The column count comes from the results, so "will it fit" cannot be settled
% while writing the caption.
\newcommand{\fittable}[1]{%
  \resizebox{\ifdim\width>\linewidth\linewidth\else\width\fi}{!}{#1}}

\newtheorem{theorem}{Theorem}
\newtheorem{lemma}{Lemma}
\newtheorem{proposition}{Proposition}
\theoremstyle{definition}
\newtheorem{definition}{Definition}
\theoremstyle{remark}
\newtheorem{remark}{Remark}

\DeclareMathOperator{\KL}{KL}
\DeclareMathOperator{\TV}{TV}
\DeclareMathOperator{\NB}{NB}
\DeclareMathOperator*{\argmax}{arg\,max}
\newcommand{\Vc}{\mathcal{V}}
\newcommand{\IV}{I_{\Vc}}
\newcommand{\HV}{H_{\Vc}}
\newcommand{\Xset}[1]{X_{#1}}
\newcommand{\pvi}{\mathrm{pvi}}
\newcommand{\bits}{\,\mathrm{bits}}


% Include a figure if the analysis that produces it has been run, otherwise leave
% a visible placeholder. This lets the manuscript compile from a partial results
% tree without ever silently dropping a panel.
\newcommand{\infogainfig}[2][\textwidth]{%
  \IfFileExists{#2.pdf}{\includegraphics[width=#1]{#2}}{%
    \IfFileExists{#2.png}{\includegraphics[width=#1]{#2}}{%
      \fbox{\parbox[c][0.26\textheight][c]{\dimexpr#1-2\fboxsep-2\fboxrule\relax}{%
        \centering\ttfamily\small
        [figure \detokenize{#2} not yet generated;
         run \detokenize{python scripts/run_all.py}]}}}}%
}

% Auto-generated by scripts/make_paper_numbers.py -- do not edit.
% Quantities from stages that have not run typeset as a visible marker
% rather than as a stale value or an undefined control sequence.
\providecommand{\todonum}[1]{\textbf{[?#1]}}

\newcommand{\thyAllHold}{yes}
\newcommand{\thyTrialsJS}{20000}
\newcommand{\thyIdentityRelErr}{1.0e-05}
\newcommand{\thyIdentityTrials}{150}
\newcommand{\thyPviErr}{1.7e-03}
\newcommand{\simRecovery}{82.9}
\newcommand{\simRecoveryCorrected}{96.0}
\newcommand{\simRegimeAccuracy}{100}
\newcommand{\simMaxN}{36{,}000}
\newcommand{\simThmOneHolds}{yes}
\newcommand{\simThmTwoHolds}{yes}
\newcommand{\simThmThreeRelErr}{3.8e-08}
\newcommand{\simRecoveryFloor}{0.01}
\newcommand{\simSubsetsScored}{11}
\newcommand{\simSubsetsBelowFloor}{0}
\newcommand{\simRecoveryFullPanel}{68.2}
\newcommand{\simMeanAbsErr}{0.0097}
\newcommand{\simRegimeMaxN}{50{,}000}
\newcommand{\simSynergyFirstN}{18{,}000}
\newcommand{\simRedundancyFirstN}{18{,}000}
\newcommand{\powBothFrac}{42}
\newcommand{\powEventsMin}{652}
\newcommand{\powEventsMax}{3569}
\newcommand{\powSynMin}{20}
\newcommand{\powSynMax}{82}
\newcommand{\powTreeSynMax}{72}
\newcommand{\restrictedShareMin}{52}
\newcommand{\restrictedShareMax}{98}
\newcommand{\powMidAN}{18{,}000}
\newcommand{\powMidAEvents}{1{,}391}
\newcommand{\powMidARecovery}{39}
\newcommand{\powMidBN}{50{,}000}
\newcommand{\powMidBEvents}{1{,}673}
\newcommand{\powMidBRecovery}{63}
\newcommand{\regimeSynSmall}{0}
\newcommand{\regimeSynSmallN}{6{,}000}
\newcommand{\regimeSynMid}{50}
\newcommand{\regimeSynMidN}{18{,}000}
\newcommand{\regimeSynLarge}{100}
\newcommand{\regimeSynLargeN}{50{,}000}
\newcommand{\acqLabs}{87}
\newcommand{\acqEcg}{67}
\newcommand{\acqCxr}{56}
\newcommand{\acqEcho}{23}
\newcommand{\thyThmOneInformative}{92.5}
\newcommand{\exGlobalNb}{0.146}
\newcommand{\exLocalNb}{0.061}
\newcommand{\tgtSignalRhoMin}{0.25}
\newcommand{\tgtSignalRhoMax}{0.50}
\newcommand{\tgtSignalCaptureMin}{65}
\newcommand{\tgtSignalCaptureMax}{74}
\newcommand{\tgtSignalChanceMin}{53}
\newcommand{\tgtSignalChanceMax}{54}
\newcommand{\tgtCeilingMax}{0.15}
\newcommand{\tgtN}{50{,}000}
\newcommand{\flipNPairs}{1}
\newcommand{\flipPairA}{laboratory panel}
\newcommand{\flipPairB}{electrocardiogram}
\newcommand{\flipSynEndpoint}{heart-failure readmission}
\newcommand{\flipRedEndpoint}{ICU transfer}
\newcommand{\flipSynBits}{0.0024}
\newcommand{\flipRedBits}{0.0004}
\newcommand{\ablSpeedup}{4}
\newcommand{\ablNCohorts}{2}
\newcommand{\switchNKeys}{691}
\newcommand{\switchMedianChange}{2.2}
\newcommand{\switchIncreased}{45}
\newcommand{\switchNSignFlips}{7}
\newcommand{\switchNMaterial}{0}
\newcommand{\switchNRegimeFlips}{4}
\newcommand{\permNPerm}{5{,}000}
\newcommand{\permFloor}{0.0002}
\newcommand{\decEchoNullPMin}{0.917}
\newcommand{\thyWitnessResidual}{6\times10^{-17}}
\newcommand{\certTightenMin}{3.2}
\newcommand{\certTightenMax}{9.9}
\newcommand{\shapleyEffResidual}{10^{-16}}
\newcommand{\redAvoidedMin}{8}
\newcommand{\redAvoidedMax}{75}
\newcommand{\redIcuTransferAvoided}{50}
\newcommand{\redIcuTransferSaved}{379}
\newcommand{\redIcuTransferCert}{0.0171}
\newcommand{\aurocIcuTransfer}{0.752}
\newcommand{\redAkiAvoided}{75}
\newcommand{\redAkiSaved}{419}
\newcommand{\redAkiCert}{0.0330}
\newcommand{\aurocAki}{0.762}
\newcommand{\redHfReadmissionAvoided}{8}
\newcommand{\redHfReadmissionSaved}{114}
\newcommand{\redHfReadmissionCert}{0.0164}
\newcommand{\aurocHfReadmission}{0.734}
\newcommand{\redMortalityAvoided}{29}
\newcommand{\redMortalitySaved}{352}
\newcommand{\redMortalityCert}{0.0160}
\newcommand{\aurocMortality}{0.798}
\newcommand{\cohortN}{50{,}000}
\newcommand{\cohortSubjects}{50{,}000}
\newcommand{\cohortSource}{synthetic}
\newcommand{\primaryPrevalence}{7.1}
\newcommand{\primaryBaselineBits}{0.0173}
\newcommand{\primaryFullBits}{0.0607}
\newcommand{\primaryBaselineAuroc}{0.665}
\newcommand{\primaryFullAuroc}{0.798}
\newcommand{\primaryNullGap}{0.00028}
\newcommand{\primaryIdentityRelErr}{3.9e-08}
\newcommand{\primaryRestrictedBits}{0.0329}
\newcommand{\primaryRestrictedShare}{94}
\newcommand{\primarySeedShare}{11}
\newcommand{\primaryMonoViolations}{13}
\newcommand{\primaryMonoWorst}{0.0009}
\newcommand{\redTestsAvoided}{29.0}
\newcommand{\redCostSaved}{351.64}
\newcommand{\redCostSavedPct}{77.3}
\newcommand{\redAuroc}{0.7961}
\newcommand{\redAurocAll}{0.7975}
\newcommand{\redCertifiedLoss}{0.01600}
\newcommand{\redInfoGap}{0.0015}
\newcommand{\gainMaxGini}{0.45}
\newcommand{\gainMinGini}{0.34}
\newcommand{\interNRedundant}{1}
\newcommand{\interNSynergistic}{1}
\newcommand{\interNPairs}{6}
\newcommand{\interMostRedundantPair}{labs--cxr}
\newcommand{\interMostRedundantBits}{0.0048}
\newcommand{\interRegimeFlips}{1}
\newcommand{\interNEndpoints}{4}
\newcommand{\tgtLabsRhoTruth}{0.25}
\newcommand{\tgtLabsRhoRetro}{0.075}
\newcommand{\tgtLabsCeiling}{0.150}
\newcommand{\tgtLabsTrueBits}{0.0382}
\newcommand{\tgtLabsCapture}{65}
\newcommand{\tgtLabsChance}{53}
\newcommand{\decLabsMarginal}{0.0265}
\newcommand{\decLabsConditional}{0.0231}
\newcommand{\decLabsRedundant}{0.0034}
\newcommand{\decLabsSynergistic}{0.0000}
\newcommand{\decLabsRedFrac}{13}
\newcommand{\gainLabsMean}{0.0243}
\newcommand{\gainLabsPNinetynine}{0.0875}
\newcommand{\gainLabsGini}{0.34}
\newcommand{\gainLabsConcentration}{4}
\newcommand{\gainLabsFracAbove}{90.3}
\newcommand{\tgtEcgRhoTruth}{0.11}
\newcommand{\tgtEcgRhoRetro}{-0.001}
\newcommand{\tgtEcgCeiling}{-0.008}
\newcommand{\tgtEcgTrueBits}{0.0002}
\newcommand{\tgtEcgCapture}{23}
\newcommand{\tgtEcgChance}{16}
\newcommand{\decEcgMarginal}{-0.0005}
\newcommand{\decEcgConditional}{0.0150}
\newcommand{\decEcgRedundant}{0.0000}
\newcommand{\decEcgSynergistic}{0.0155}
\newcommand{\decEcgRedFrac}{0}
\newcommand{\gainEcgMean}{0.0035}
\newcommand{\gainEcgPNinetynine}{0.0203}
\newcommand{\gainEcgGini}{0.45}
\newcommand{\gainEcgConcentration}{6}
\newcommand{\gainEcgFracAbove}{4.9}
\newcommand{\tgtCxrRhoTruth}{0.50}
\newcommand{\tgtCxrRhoRetro}{0.073}
\newcommand{\tgtCxrCeiling}{0.096}
\newcommand{\tgtCxrTrueBits}{0.0210}
\newcommand{\tgtCxrCapture}{74}
\newcommand{\tgtCxrChance}{54}
\newcommand{\decCxrMarginal}{0.0069}
\newcommand{\decCxrConditional}{0.0171}
\newcommand{\decCxrRedundant}{0.0000}
\newcommand{\decCxrSynergistic}{0.0103}
\newcommand{\decCxrRedFrac}{0}
\newcommand{\gainCxrMean}{0.0101}
\newcommand{\gainCxrPNinetynine}{0.0417}
\newcommand{\gainCxrGini}{0.44}
\newcommand{\gainCxrConcentration}{4}
\newcommand{\gainCxrFracAbove}{36.2}
\newcommand{\tgtEchoRhoTruth}{0.08}
\newcommand{\tgtEchoRhoRetro}{0.035}
\newcommand{\tgtEchoCeiling}{-0.001}
\newcommand{\tgtEchoTrueBits}{0.0002}
\newcommand{\tgtEchoCapture}{19}
\newcommand{\tgtEchoChance}{16}
\newcommand{\decEchoMarginal}{-0.0002}
\newcommand{\decEchoConditional}{-0.0005}
\newcommand{\decEchoRedundant}{0.0003}
\newcommand{\decEchoSynergistic}{0.0000}
\newcommand{\decEchoRedFrac}{176}
\newcommand{\gainEchoMean}{0.0030}
\newcommand{\gainEchoPNinetynine}{0.0146}
\newcommand{\gainEchoGini}{0.45}
\newcommand{\gainEchoConcentration}{5}
\newcommand{\gainEchoFracAbove}{3.9}
\newcommand{\tgtNotesRhoTruth}{\todonum{tgtNotesRhoTruth}}
\newcommand{\tgtNotesRhoRetro}{\todonum{tgtNotesRhoRetro}}
\newcommand{\tgtNotesCeiling}{\todonum{tgtNotesCeiling}}
\newcommand{\tgtNotesTrueBits}{\todonum{tgtNotesTrueBits}}
\newcommand{\tgtNotesCapture}{\todonum{tgtNotesCapture}}
\newcommand{\tgtNotesChance}{\todonum{tgtNotesChance}}
\newcommand{\decNotesMarginal}{\todonum{decNotesMarginal}}
\newcommand{\decNotesConditional}{\todonum{decNotesConditional}}
\newcommand{\decNotesRedundant}{\todonum{decNotesRedundant}}
\newcommand{\decNotesSynergistic}{\todonum{decNotesSynergistic}}
\newcommand{\decNotesRedFrac}{\todonum{decNotesRedFrac}}
\newcommand{\gainNotesMean}{\todonum{gainNotesMean}}
\newcommand{\gainNotesPNinetynine}{\todonum{gainNotesPNinetynine}}
\newcommand{\gainNotesGini}{\todonum{gainNotesGini}}
\newcommand{\gainNotesConcentration}{\todonum{gainNotesConcentration}}
\newcommand{\gainNotesFracAbove}{\todonum{gainNotesFracAbove}}
\newcommand{\ablAttnFmSyn}{58}
\newcommand{\ablAttnFmSynLo}{56}
\newcommand{\ablAttnFmSynHi}{59}
\newcommand{\ablAttnFmRegimes}{8/8}
\newcommand{\ablAttnFmSecs}{1{,}259}
\newcommand{\ablAttnNoFmSyn}{6}
\newcommand{\ablAttnNoFmSynLo}{6}
\newcommand{\ablAttnNoFmSynHi}{6}
\newcommand{\ablAttnNoFmRegimes}{4/8}
\newcommand{\ablAttnNoFmSecs}{870}
\newcommand{\ablConcatFmSyn}{57}
\newcommand{\ablConcatFmSynLo}{48}
\newcommand{\ablConcatFmSynHi}{66}
\newcommand{\ablConcatFmRegimes}{8/8}
\newcommand{\ablConcatFmSecs}{309}
\newcommand{\ablConcatNoFmSyn}{46}
\newcommand{\ablConcatNoFmSynLo}{39}
\newcommand{\ablConcatNoFmSynHi}{53}
\newcommand{\ablConcatNoFmRegimes}{8/8}
\newcommand{\ablConcatNoFmSecs}{334}
\newcommand{\decLabsMortalityMarg}{0.0265}
\newcommand{\decLabsMortalityCond}{0.0231}
\newcommand{\decLabsMortalityRed}{0.0034}
\newcommand{\decLabsMortalitySyn}{0.0000}
\newcommand{\decLabsMortalityP}{0.0002}
\newcommand{\decEcgMortalityMarg}{-0.0005}
\newcommand{\decEcgMortalityCond}{0.0150}
\newcommand{\decEcgMortalityRed}{0.0000}
\newcommand{\decEcgMortalitySyn}{0.0155}
\newcommand{\decEcgMortalityP}{0.0002}
\newcommand{\decCxrMortalityMarg}{0.0069}
\newcommand{\decCxrMortalityCond}{0.0171}
\newcommand{\decCxrMortalityRed}{0.0000}
\newcommand{\decCxrMortalitySyn}{0.0103}
\newcommand{\decCxrMortalityP}{0.0002}
\newcommand{\decEchoMortalityMarg}{-0.0002}
\newcommand{\decEchoMortalityCond}{-0.0005}
\newcommand{\decEchoMortalityRed}{0.0003}
\newcommand{\decEchoMortalitySyn}{0.0000}
\newcommand{\decEchoMortalityP}{0.9946}
\newcommand{\decNotesMortalityMarg}{\todonum{decNotesMortalityMarg}}
\newcommand{\decNotesMortalityCond}{\todonum{decNotesMortalityCond}}
\newcommand{\decNotesMortalityRed}{\todonum{decNotesMortalityRed}}
\newcommand{\decNotesMortalitySyn}{\todonum{decNotesMortalitySyn}}
\newcommand{\decNotesMortalityP}{\todonum{decNotesMortalityP}}
\newcommand{\decLabsHfReadmissionMarg}{0.0502}
\newcommand{\decLabsHfReadmissionCond}{0.0242}
\newcommand{\decLabsHfReadmissionRed}{0.0261}
\newcommand{\decLabsHfReadmissionSyn}{0.0000}
\newcommand{\decLabsHfReadmissionP}{0.0002}
\newcommand{\decEcgHfReadmissionMarg}{0.0103}
\newcommand{\decEcgHfReadmissionCond}{0.0156}
\newcommand{\decEcgHfReadmissionRed}{0.0000}
\newcommand{\decEcgHfReadmissionSyn}{0.0054}
\newcommand{\decEcgHfReadmissionP}{0.0002}
\newcommand{\decCxrHfReadmissionMarg}{0.0312}
\newcommand{\decCxrHfReadmissionCond}{0.0039}
\newcommand{\decCxrHfReadmissionRed}{0.0274}
\newcommand{\decCxrHfReadmissionSyn}{0.0000}
\newcommand{\decCxrHfReadmissionP}{0.0002}
\newcommand{\decEchoHfReadmissionMarg}{0.0090}
\newcommand{\decEchoHfReadmissionCond}{0.0126}
\newcommand{\decEchoHfReadmissionRed}{0.0000}
\newcommand{\decEchoHfReadmissionSyn}{0.0036}
\newcommand{\decEchoHfReadmissionP}{0.0002}
\newcommand{\decNotesHfReadmissionMarg}{\todonum{decNotesHfReadmissionMarg}}
\newcommand{\decNotesHfReadmissionCond}{\todonum{decNotesHfReadmissionCond}}
\newcommand{\decNotesHfReadmissionRed}{\todonum{decNotesHfReadmissionRed}}
\newcommand{\decNotesHfReadmissionSyn}{\todonum{decNotesHfReadmissionSyn}}
\newcommand{\decNotesHfReadmissionP}{\todonum{decNotesHfReadmissionP}}
\newcommand{\decLabsAkiMarg}{0.1215}
\newcommand{\decLabsAkiCond}{0.1143}
\newcommand{\decLabsAkiRed}{0.0072}
\newcommand{\decLabsAkiSyn}{0.0000}
\newcommand{\decLabsAkiP}{0.0002}
\newcommand{\decEcgAkiMarg}{-0.0004}
\newcommand{\decEcgAkiCond}{-0.0009}
\newcommand{\decEcgAkiRed}{0.0005}
\newcommand{\decEcgAkiSyn}{0.0000}
\newcommand{\decEcgAkiP}{0.9994}
\newcommand{\decCxrAkiMarg}{0.0091}
\newcommand{\decCxrAkiCond}{0.0014}
\newcommand{\decCxrAkiRed}{0.0078}
\newcommand{\decCxrAkiSyn}{0.0000}
\newcommand{\decCxrAkiP}{0.0036}
\newcommand{\decEchoAkiMarg}{-0.0006}
\newcommand{\decEchoAkiCond}{-0.0009}
\newcommand{\decEchoAkiRed}{0.0003}
\newcommand{\decEchoAkiSyn}{0.0000}
\newcommand{\decEchoAkiP}{1.0000}
\newcommand{\decNotesAkiMarg}{\todonum{decNotesAkiMarg}}
\newcommand{\decNotesAkiCond}{\todonum{decNotesAkiCond}}
\newcommand{\decNotesAkiRed}{\todonum{decNotesAkiRed}}
\newcommand{\decNotesAkiSyn}{\todonum{decNotesAkiSyn}}
\newcommand{\decNotesAkiP}{\todonum{decNotesAkiP}}
\newcommand{\decLabsIcuTransferMarg}{0.0149}
\newcommand{\decLabsIcuTransferCond}{0.0150}
\newcommand{\decLabsIcuTransferRed}{0.0000}
\newcommand{\decLabsIcuTransferSyn}{0.0001}
\newcommand{\decLabsIcuTransferP}{0.0002}
\newcommand{\decEcgIcuTransferMarg}{-0.0004}
\newcommand{\decEcgIcuTransferCond}{0.0007}
\newcommand{\decEcgIcuTransferRed}{0.0000}
\newcommand{\decEcgIcuTransferSyn}{0.0011}
\newcommand{\decEcgIcuTransferP}{0.0010}
\newcommand{\decCxrIcuTransferMarg}{0.0189}
\newcommand{\decCxrIcuTransferCond}{0.0208}
\newcommand{\decCxrIcuTransferRed}{0.0000}
\newcommand{\decCxrIcuTransferSyn}{0.0019}
\newcommand{\decCxrIcuTransferP}{0.0002}
\newcommand{\decEchoIcuTransferMarg}{-0.0001}
\newcommand{\decEchoIcuTransferCond}{-0.0002}
\newcommand{\decEchoIcuTransferRed}{0.0001}
\newcommand{\decEchoIcuTransferSyn}{0.0000}
\newcommand{\decEchoIcuTransferP}{0.9168}
\newcommand{\decNotesIcuTransferMarg}{\todonum{decNotesIcuTransferMarg}}
\newcommand{\decNotesIcuTransferCond}{\todonum{decNotesIcuTransferCond}}
\newcommand{\decNotesIcuTransferRed}{\todonum{decNotesIcuTransferRed}}
\newcommand{\decNotesIcuTransferSyn}{\todonum{decNotesIcuTransferSyn}}
\newcommand{\decNotesIcuTransferP}{\todonum{decNotesIcuTransferP}}


\title{\bfseries INFOGAIN: decomposing the information value of multimodal
diagnostic testing at the level of the individual patient}

\author[1]{Author One}
\author[1]{Author Two}
\affil[1]{Affiliation}

\date{}

\begin{document}
\maketitle

\begin{abstract}
\noindent
Clinical medicine orders tests in bundles, but the value of a test is not a
property of the test: it depends on what is already known about the patient in
front of you. We formalise that dependence. Using $\Vc$-usable information --
the information a bounded predictor can actually extract -- we define a
decomposition of a diagnostic modality's value, \emph{relative to the tests
already in hand}, into a redundant part the context already contained and a
synergistic part that exists only in combination. Unlike partial information
decomposition, every term is identified by a decision problem rather than by a
contested axiom, and every term is estimable from data. We then prove that these
bits are clinical currency, not a metaphor. Theorem~3 is an identity: the
information a test adds equals the integral of its net-benefit gain over all
risk thresholds a clinician might hold, $\Delta I = \int_0^1 \Delta \NB(t)\,
\mathrm{d}t/t$. Theorem~1 turns synergy into a certified lower bound on
achievable discrimination; Theorem~2 turns redundancy into a certificate that
skipping a test cannot cost more than $\sqrt{\varepsilon/2}/(1-t)$ in net
benefit. We estimate all quantities patient-by-patient with a cross-fitted
masked multimodal model, correct the finite-sample learning bias by
learning-curve extrapolation, and validate the whole pipeline against a
generative model in which the true decomposition is known in closed form -- which
also shows what governs whether synergy is visible at all: not the cohort size
but the number of events for which both modalities were acquired, a quantity no
multimodal study reports. Applied to a cohort of \cohortN\ encounters with
electrocardiograms, chest radiographs, laboratory panels and notes for the same
patients, across four endpoints, the same pair of modalities is redundant for
one and synergistic for another, and the benefit of any given test is
concentrated in a minority of patients (Gini \gainMaxGini). Ordering tests by
patient-level expected information gain reaches the discrimination of testing
everyone while ordering \redTestsAvoided\% fewer tests, with a certified
net-benefit loss below \redCertifiedLoss. The reported cohort is the calibrated
simulator, not MIMIC-IV: the extraction targets linked MIMIC-IV records and runs
through the identical code path, but requires credentialed access we do not yet
hold (see Data status). The framework converts ``which test should this patient
get next'' from a habit into a computation.
\end{abstract}

\paragraph{Data status.}
\emph{The numbers in this manuscript were produced by the analysis pipeline in
this repository run on the calibrated simulator described in
Methods~\ref{sec:sim}, which reproduces the MIMIC-IV cohort's schema, modality
dimensions, endpoint prevalences and informative test-ordering, and for which
the true information decomposition is available in closed form. The MIMIC-IV
extraction code is complete and configured
(\texttt{infogain.data.mimic\_cohort}); running it requires credentialed
PhysioNet access, and the corresponding numbers will replace these before
submission. Every table, figure and theorem check is produced by the identical
code path on either input.}

% ===================================================================== %
\section{Introduction}

A patient arrives in the emergency department. Their vital signs are already
recorded; an electrocardiogram takes minutes and costs little; a chest
radiograph costs more and irradiates; an echocardiogram costs an order of
magnitude more and takes hours. Which should be ordered? The question is asked
millions of times a day and answered almost entirely by habit, protocol and
availability. Between a fifth and a third of health spending in high-income
systems is estimated to buy nothing \citep{oecd2017tackling,brownlee2017overuse},
and diagnostic testing is a large and growing share of it
\citep{cassel2012choosing}. At the same time, the opposite failure -- a test not
ordered for a patient who needed it -- is the one that reaches the coroner.

Machine learning has approached this from the wrong end. The dominant paradigm
in multimodal clinical AI is to fuse everything available and report that the
fused model beats each single modality
\citep{huang2020fusion,soenksen2022holisticai,acosta2022multimodal}. That
comparison answers a question nobody asks. It tells us that more data helps on
average; it does not tell us \emph{which} patient needed \emph{which} test, and
it cannot, because a cohort-average performance difference is not a
patient-level decision rule. It also systematically overstates the case: a
modality that improves mean AUROC by 0.01 might be doing so by helping 3\% of
patients enormously and the rest not at all -- or by helping everyone
imperceptibly. Those two worlds imply opposite testing policies and are
indistinguishable in the reported statistic.

What is needed is a quantity that is (i) defined per patient, (ii) defined
\emph{conditionally} on what has already been measured, and (iii) convertible
into the units a clinician or a payer actually optimises. Information theory
supplies a natural candidate, and the partial information decomposition
literature \citep{williams2010pid,bertschinger2014uniqueinfo,griffith2014synergy}
supplies exactly the vocabulary -- unique, redundant, synergistic -- that the
clinical question calls for. It also supplies an unresolved twenty-year dispute
about which redundancy functional is correct, which makes any PID-based clinical
claim vulnerable to the objection that a different, equally axiomatic choice
would have given a different answer.

We take a different route on all three counts.

\textbf{Computationally bounded information.} We work with $\Vc$-usable
information \citep{xu2020vinfo,ethayarajh2022pvi} rather than Shannon mutual
information. Mutual information asks what an unbounded observer could extract; a
hospital deploys a specific model on a specific sample size, and the bits it
cannot reach are not clinically real. $\Vc$-information is defined relative to a
predictor family, has a pointwise form that assigns a value to each individual
patient, and -- crucially -- is a lower bound by construction, so an
under-trained model understates rather than overstates the value of a test.

\textbf{Context-relative decomposition.} We do not attempt to define redundancy
in the abstract. We define it \emph{relative to a context} $S$ of tests already
performed: the redundant part of a modality's value is the part its marginal
value loses once $S$ is conditioned on, and the synergistic part is the part
that only appears then. Both are differences of estimable quantities, both are
non-negative, at most one is non-zero (they coincide at zero exactly when the
context changes nothing), and together they reconstruct the conditional value
exactly. Nothing is left to axiom choice. Shapley averaging
over acquisition orders \citep{shapley1953value,lundberg2017shap} removes the
remaining arbitrariness of which context to condition on.

\textbf{Bits as clinical currency.} The step that makes information theory
clinically actionable is an identity, not an analogy. Writing $\NB(t)$ for
Vickers' net benefit at risk threshold $t$ \citep{vickers2006decision}, we show
(Theorem~\ref{thm:identity}) that the information a test adds is exactly the
integral of the net-benefit it adds, weighted across every threshold a clinician
might hold:
\[
\Delta I \;=\; \int_0^1 \Delta \NB(t)\,\frac{\mathrm{d}t}{t}.
\]
This follows from Schervish's representation of proper scoring rules
\citep{schervish1989general} and it has three consequences we exploit. It says
which bits matter -- those earned at thresholds a clinician would actually act
on, defining a \emph{clinically restricted usable information}. It converts a
synergy estimate into a certified lower bound on achievable discrimination
(Theorem~\ref{thm:auroc}). And it converts a redundancy estimate into a
certificate that omitting a test cannot cost more than a stated amount of net
benefit (Theorem~\ref{thm:omit}) -- the safety guarantee that any
``order fewer tests'' claim requires and almost never provides.

\textbf{Per-patient, prospectively.} Finally, the score has to be computable
\emph{before} the test is done, without its result and without the outcome. We
estimate the expected information gain
$\Delta_i(m\mid S) = \mathbb{E}_{x_m}\KL(f_{S\cup m}\|f_S)$ by integrating over
what the test would plausibly have shown, resampled from real results of similar
patients who did receive it (Methods~\ref{sec:gain}). Averaged over patients it
must return the cohort conditional gain -- an identity we verify rather than
assume, and one that jointly tests the sampler, the model family and the
cross-fitting.

MIMIC-IV is the setting that makes this possible. It is, to our knowledge, the
only public resource in which electrocardiograms \citep{gow2023mimicecg}, chest
radiographs \citep{johnson2019mimiccxr}, laboratory and medication records
\citep{johnson2023mimiciv} and free-text notes \citep{johnson2023mimicnote} are
all linked to the same patients on a common clock. Without that linkage the
conditional question -- what does the radiograph add \emph{given} the
electrocardiogram, \emph{for this patient} -- cannot be posed at all.

% ===================================================================== %
\section{Results}

\subsection{A decomposition whose every term is a decision}
\label{sec:decomp}

Let $Y$ be a binary endpoint, $M$ the set of available modalities, and $\Vc$ a
family of predictors of $Y$ from any subset of $M$. The $\Vc$-usable information
in a subset $S$ is $\IV(\Xset{S}\to Y) = \HV(Y) - \HV(Y\mid \Xset{S})$, each term
an infimum of expected log-loss over $\Vc$ \citep{xu2020vinfo}. We realise $\Vc$
as a single network with a learned ``absent'' token per modality
(Methods~\ref{sec:family}), so that $f_S$ exists for every $S$ and the family is
closed under masking; closure is what makes $\IV$ monotone on the subset lattice
and the conditional gains below non-negative in population.

Fix a free baseline $S_0$ (demographics and vital signs: acquired for everyone,
never a decision) and a context $S \supseteq S_0$ of tests already performed. For
a candidate modality $m$ we define
\begin{align}
I_m &:= \IV(\Xset{S_0\cup m}\to Y) - \IV(\Xset{S_0}\to Y)
  &&\text{(marginal: value alone)},\\
U(m\mid S) &:= \IV(\Xset{S\cup m}\to Y) - \IV(\Xset{S}\to Y)
  &&\text{(conditional: value in context)},\\
R(m\mid S) &:= \big(I_m - U(m\mid S)\big)_+ &&\text{(redundant)},\\
\mathrm{Syn}(m\mid S) &:= \big(U(m\mid S) - I_m\big)_+ &&\text{(synergistic)},
\end{align}
so that $U(m\mid S) = I_m - R(m\mid S) + \mathrm{Syn}(m\mid S)$ identically, with
$R\cdot\mathrm{Syn}=0$. Each term answers a decision: $I_m$ is what a
single-modality study would report; $U$ is what the test is worth to a clinician
who already has $S$; $R$ is what they would be paying for twice;
$\mathrm{Syn}$ is what they would miss by evaluating the test in isolation.

Because $U(m\mid S)$ depends on the accident of which tests happened to be done
first, we also report the Shapley value $\varphi_m$ of each modality over all
acquisition orders. Efficiency guarantees $\sum_m \varphi_m$ equals the total
usable information of the full panel, so no modality can be credited with bits
the panel does not contain; in our runs the efficiency residual is numerically
zero ($<\shapleyEffResidual\bits$).

\subsection{Information is exactly integrated net benefit}
\label{sec:res-identity}

The bridge from bits to bedside is an identity. Let $\eta_1 = P(Y=1\mid X_S)$ and
$\eta_2 = P(Y=1\mid X_{S\cup m})$, and let $\NB_j(t) = \mathbb{E}[(\eta_j-t)_+]/(1-t)$
be the net benefit of acting at risk threshold $t$.

\begin{theorem}[Information--net-benefit identity]
\label{thm:identity}
If $\eta_1 = \mathbb{E}[\eta_2\mid X_S]$, then in nats
\[
\Delta I \;=\; \mathbb{E}\,\KL\!\big(\mathrm{Ber}(\eta_2)\,\|\,\mathrm{Ber}(\eta_1)\big)
\;=\; \int_0^1 \big[\NB_2(t)-\NB_1(t)\big]\,\frac{\mathrm{d}t}{t}.
\]
\end{theorem}

The proof (Supplementary~\S\ref{sec:s-identity}) applies Schervish's representation of the negative
Shannon entropy as a mixture of threshold losses, then uses the nesting condition
to kill the linear term. Numerically the identity holds to a relative error of
\thyIdentityRelErr\ over \thyIdentityTrials\ random nested posterior pairs, and
to \primaryIdentityRelErr\ on the fitted cohort models
(Fig.~\ref{fig:identity}).

Theorem~\ref{thm:identity} immediately says which bits are worth counting. The
weight $1/t$ concentrates on small thresholds, and a clinician holds a threshold
in a bounded range -- roughly 2--30\% for the decisions studied here. We
therefore report the \emph{clinically restricted usable information}
$I^{[a,b]} = \int_a^b \Delta\NB(t)\,\mathrm{d}t/t$.

How much this matters depends entirely on the endpoint, which is itself the
argument for computing it rather than assuming it. For 30-day mortality
\primaryRestrictedShare\% of the full panel's information gain is redeemable
inside $[0.02,0.30]$, so the restricted and unrestricted quantities nearly
coincide; for acute kidney injury only \restrictedShareMin\% is, and the
remaining half is earned at thresholds no clinician acts on. Across the four
endpoints the redeemable share runs from \restrictedShareMin\% to
\restrictedShareMax\%. An unrestricted bit count is therefore not a fixed
overstatement that a constant could correct: it overstates clinical value by an
endpoint-specific factor, and nothing in the number itself says which case you
are in.

\subsection{Synergy bounds achievable discrimination; redundancy bounds safe omission}

AUROC is not a proper scoring rule, so no unconditional map from bits to AUROC
gain can exist -- a model can gain information by recalibrating and leave every
ranking untouched. Stating otherwise is a common error. What is true, and
sufficient, is a stratified statement.

\begin{theorem}[Synergy $\Rightarrow$ achievable discrimination]
\label{thm:auroc}
Let $A^*_T$ be the maximal AUROC achievable from $X_T$. For any partition
$\{B_k\}$ of the range of $\eta_S$ into intervals,
\[
A^*_{S\cup m} \;\ge\; A(\eta_S) \;+\; \sum_k w_k\big[A_k(\eta_{S\cup m}) - A_k(\eta_S)\big],
\]
where $w_k$ is the share of case--control pairs inside $B_k$ and $A_k$ the
within-bin AUROC. Within each bin the second term is bounded below by
$U_k / \big(2H(\pi_k)\big)$, the within-bin conditional usable information.
\end{theorem}

The proof is a ranking argument (Supplementary~\S\ref{sec:s-lemmas}): the lexicographic score
``bin index, ties broken by $\eta_{S\cup m}$'' is $X_{S\cup m}$-measurable and
agrees with $\eta_S$ on every across-bin pair, so it can only differ within bins.
The information step uses two lemmas we state and verify separately -- a
Kolmogorov--Smirnov sandwich $D^+ \le 2A-1 \le 2D^+$ (the lower bound for
concave ROCs, hence for the ROC convex hull \citep{provost2001robust}) and a
Jensen--Shannon bound $8\pi^2(1-\pi)^2\TV^2 \le I \le H(\pi)\TV$
\citep{lin1991divergence}. Both directions of the latter have short proofs
(Supplementary~\S\ref{sec:s-lemmas}) resting on $\mathbb{E}[\eta]=\pi$: the upper bound from the
two chords of the concave binary entropy, the lower from Pinsker plus Jensen. A
sharper lower constant held in every numerical check but has no proof, so we do
not use it; only the upper bound enters Theorem~\ref{thm:auroc}, and a looser
lower bound merely widens the (unused) upper end of the bracket.

The theorem's own content is an \emph{identity}: the certified sum equals the
AUROC gain of an explicit witness score, the lexicographic
``$\eta_S$-bin, ties broken by $\eta_{S\cup m}$''. We verify that identity
rather than the inequality it implies, and it holds to $\thyWitnessResidual$. All
four lemma inequalities hold with zero violations over $\thyTrialsJS$
Monte-Carlo trials each (Supplementary Table~\ref{tab:s-theory}), and on exact simulator
posteriors the bound never exceeded the achievable gain (\simThmOneHolds).

The converse direction is what licenses omitting a test.

\begin{theorem}[Redundancy $\Rightarrow$ certified safe omission]
\label{thm:omit}
If omitting modality $m$ forgoes at most $\varepsilon$ nats of information, then
at every risk threshold $t$ the net-benefit loss obeys
$\Delta\NB(t) \le \sqrt{\varepsilon/2}\,/\,(1-t)$. Integrating the true
threshold weight rather than its lower bound gives a strictly tighter,
threshold-specific certificate (Supplementary~\S\ref{sec:s-omit}), which on our runs is
\certTightenMin--\certTightenMax$\times$ smaller at the thresholds used in
triage ($t \le 0.05$).
\end{theorem}

The practical content is a guarantee in the direction that matters. A test whose
forgone information is small \emph{cannot} be costing much benefit, whatever the
sample happens to show; and inverting the bound gives the minimum information
gain a test must deliver before it could possibly move the decision curve by a
clinically visible amount. Both bounds dominated the realised loss in every
simulation trial (\simThmTwoHolds).

\subsection{The estimator recovers information it cannot see}
\label{sec:res-validation}

None of the decomposition terms is observable in real data, so we validated
against a generative model where all of them are available in closed form
(Methods~\ref{sec:sim}): binary latent factors drive Gaussian modality read-outs
and a logistic outcome, so posteriors for every subset follow by enumeration. The
design deliberately contains all three regimes at once -- a factor loaded on two
modalities (redundant), factors loaded on one (unique), and a centred
interaction between factors visible to different modalities (pure synergy, zero
marginal effect by construction).

Two findings matter. First, \textbf{the estimator is consistent but approaches
from below}: at $n=\simMaxN$ it recovers a median \simRecovery\% of the true
information, and the shortfall is concentrated in the synergistic terms, which
require an interaction to be learned rather than a main effect (the full panel,
which carries all of it, sits at \simRecoveryFullPanel\%). This is the expected
behaviour of a computationally bounded quantity, not a defect -- $\IV$ is an
infimum over what a finite training set can reach, so every estimate is a valid
lower bound. Learning-curve extrapolation
$\hat I(m) = I_\infty - a m^{-\alpha}$ closes much of the remaining gap, to
\simRecoveryCorrected\% (Fig.~\ref{fig:recovery}).

Recovery percentages are reported over the \simSubsetsScored\ subsets whose true
information exceeds \simRecoveryFloor\ bits (\simSubsetsBelowFloor\ fall below).
Below that floor the ratio divides one noise term by another and swings in both
directions -- a subset genuinely worth a few thousandths of a bit can show ``recovery''
far above 100\% or below 0\% purely by chance -- which says nothing about the
estimator. The mean absolute error
over \emph{all} subsets, which is well behaved at zero, is
\simMeanAbsErr\ bits.

Second, \textbf{the synergistic term is the last part of the decomposition to
become estimable, and what governs it is not the cohort size but the number of
events for which \emph{both} modalities were acquired}. An interaction can only
be learned from patients who had both tests. With the acquisition rates in our
design -- \acqEcg\% for the electrocardiogram, \acqCxr\% for the radiograph --
that overlap
is \powBothFrac\% of the cohort, and the event count inside it correspondingly
smaller.

Crossing cohort size with an all-acquired variant separates the two
explanations, and it is the overlap that tracks (Supplementary Table~\ref{tab:s-power}).
Recovery of the synergistic terms rises from \powSynMin\% at \powEventsMin\
overlapping events to \powSynMax\% at \powEventsMax. The decisive comparison is
the middle pair: a cohort of \powMidAN\ with complete acquisition
(\powMidAEvents\ overlapping events, \powMidARecovery\% recovered) and one of
\powMidBN\ with realistic acquisition (\powMidBEvents\ events,
\powMidBRecovery\% recovered) behave alike despite differing almost threefold
in $n$. What the estimator sees is the overlap.

Two things follow. For anyone designing a multimodal study, the quantity to
power on is the overlap, not the cohort: at these acquisition rates, detecting a
given interaction needs roughly two and a half times the patients the same
interaction would need if every test were universal, and no headline cohort size
reveals that. For anyone reading one, an estimate of zero synergy from a
modestly-sized cohort is weak evidence of absence -- in our own runs the
synergistic regimes are called correctly at $n=\simRegimeMaxN$ and not at
$n=\powMidAN$
(Supplementary Table~\ref{tab:s-firstcorrect}). The error is at least in the safe direction: it
understates a test's value, so it cannot manufacture a reduction that is not
there. The learning-curve extrapolation, which lifts median recovery from
\simRecovery\% to \simRecoveryCorrected\%, and the tree reference of
Supplementary Table~\ref{tab:s-tree} give two independent ways to tell "hard to reach" from
"not present".

The fusion head used throughout carries an explicit second-order term over the
modality tokens ($\sum_{m<m'} e_m \odot e_{m'}$, computable in $O(Md)$ by the
factorization-machine identity). An earlier version of this manuscript declined
to attribute its effect, because it had shipped alongside a change to the
training-mask distribution. Isolating it (Methods~\ref{sec:family},
Supplementary Table~\ref{tab:s-ablation}) gives a more specific answer than we
expected: the term is indispensable to the transformer head and does nothing
for a concatenation head, which recovers the synergistic terms perfectly well
without it. It is not that interactions need an explicit multiplicative path in
general; it is that attention routes all cross-modal coupling through one
scalar per head, where concatenation leaves the features adjacent for the head
to multiply. This bears on the estimator's
tightness rather than on any result: recovery is the same for both heads once
the term is present, and \emph{no} architecture recovers more than
\ablConcatFmSynHi\% of the synergy at this overlap, which is the point of the
preceding paragraphs.

We report three further diagnostics with every run, because each is a way the
analysis could be quietly wrong. The null-model gap -- how far the network's
all-absent member sits from the base rate -- is \primaryNullGap\ bits, which
bounds the inflation of every level estimate. Monotonicity of the estimated
lattice is violated \primaryMonoViolations\ times, worst case
\primaryMonoWorst\ bits, and we report the Euclidean projection onto the monotone
cone alongside the raw values rather than silently clipping. Between-seed
variance accounts for \primarySeedShare\% of total estimator variance and is
folded into every confidence interval; treating a single training run as the
answer would understate uncertainty by that amount.

\subsection{The same two tests are redundant for one disease and synergistic for another}

Figure~\ref{fig:decomp} gives the decomposition for the primary endpoint,
Figure~\ref{fig:map} the pairwise interaction map, and
Table~\ref{tab:decomposition} all four endpoints at once. The structure is not a
property of the modalities:

\begin{itemize}
\item For \textbf{30-day mortality} the electrocardiogram is worth nothing on
  its own (\decEcgMarginal\ bits, indistinguishable from zero) yet contributes
  \decEcgConditional\ bits once every other test is in hand -- \decEcgSynergistic\
  bits of pure synergy ($p \le \permFloor$, the smallest value a
  \permNPerm-draw sign-flip permutation test can return). A single-modality
  study would have discarded
  it. The laboratory panel is the reverse: the largest marginal value
  (\decLabsMarginal\ bits) but \decLabsRedFrac\% of it already present in the
  rest of the panel.
\item The \textbf{same modality changes role between endpoints}. The chest
  radiograph is largely redundant for acute kidney injury -- \decCxrAkiMarg\
  bits alone, only \decCxrAkiCond\ once the laboratory panel is in hand -- and
  synergistic for mortality (\decCxrMortalitySyn\ bits). The echocardiogram is
  indistinguishable from useless for mortality, AKI and ICU transfer
  ($p \ge \decEchoNullPMin$ in each) and carries \decEchoHfReadmissionCond\ bits for heart-failure
  readmission, more than it carries alone (\decEchoHfReadmissionMarg), so that
  \decEchoHfReadmissionSyn\ bits of it appear only in the presence of the other
  tests. Which test it needs we cannot say: no pairwise interaction involving the
  echocardiogram is separated from zero at this cohort size.
\item The heading's claim is literal, and \flipNPairs\ pair meets it: the
  \flipPairA\ and the \flipPairB\ are \emph{synergistic} for
  \flipSynEndpoint\ (\flipSynBits\ bits of interaction, interval clear of
  zero) and \emph{redundant} for \flipRedEndpoint\ (\flipRedBits\ bits of
  overlap). Every other pair is called for at most one endpoint and
  indeterminate elsewhere, which is the overlap constraint of the previous
  section showing up in the pairwise map: pairwise interactions are estimated
  from the patients who had both tests, and at this cohort size most pairs do
  not have enough of them to be called at all.
\item Consequently \textbf{how much testing can be avoided is a property of the
  question, not of the hospital}: at matched discrimination the reduction runs
  from \redHfReadmissionAvoided\% for heart-failure readmission, where almost
  every test carries something the others do not, to \redAkiAvoided\% for acute
  kidney injury, where the laboratory panel carries nearly all of it.
\end{itemize}

This is the empirical case against test panels defined per department rather
than per question. A bundle justified by its members' marginal performance on
one endpoint can be almost entirely redundant for the endpoint it is actually
being used to inform.

\subsection{The benefit of a test is concentrated in a minority of patients}

Figure~\ref{fig:gains} shows the distribution of the per-patient prospective
gain $\Delta_i(m\mid S_0)$. It is right-skewed for every modality, with Gini
coefficients from \gainMinGini\ to \gainMaxGini. The sharpest way to put it is in
counts rather than moments: an electrocardiogram would move the risk estimate by
more than 0.01 bits for only \gainEcgFracAbove\% of patients and an
echocardiogram for \gainEchoFracAbove\%, against \gainLabsFracAbove\% for the
laboratory panel. The electrocardiogram's mean gain is \gainEcgMean\ bits and its
99th percentile \gainEcgPNinetynine, a factor of
$\approx$\,\gainEcgConcentration.

This concentration is what makes patient-level targeting worth doing at all: had
the gain been uniform, no policy could beat a blanket one and the cohort average
would have been the whole story. It is also not simply risk stratification.
Figure~\ref{fig:byrisk} plots the mean gain against baseline risk decile: the
curves are neither flat nor identically shaped, so ordering by predicted risk --
the standard clinical heuristic -- is a different, and worse, policy than
ordering by expected information gain.

The prospective score, computed without the test result or the outcome, agrees
with the retrospective realised gain in cohort mean, as the tower property
requires (Supplementary Table~\ref{tab:s-consistency}); the two are computed by entirely different
routes and their agreement jointly validates the conditional sampler, the model
family and the cross-fitting.

Their \emph{rank} correlation, in the same table, is close to zero, and the
natural reading -- that the score is noise patient by patient -- is wrong for a
reason worth stating, because the same table will appear in any paper that
attempts this check. The retrospective quantity
$\mathrm{pvi}_i(S\cup m)-\mathrm{pvi}_i(S)$ is a single-draw log-likelihood
ratio evaluated at the one outcome this patient happened to have. Its
conditional expectation given $x_i$ is the gain, which is why the means agree,
but for a rare endpoint the draw swamps the signal and it cannot rank patients
however good the score is. The simulator settles this, because there the exact
per-patient gain is available: at $n=\tgtN$ the correlation the \emph{exact}
gain achieves against the retrospective quantity is at most \tgtCeilingMax\ --
\tgtLabsCeiling\ for the laboratory panel, \tgtCxrCeiling\ for the radiograph,
indistinguishable from zero for the other two. Our estimator scores
\tgtLabsRhoRetro\ and \tgtCxrRhoRetro\ against the same yardstick
(Supplementary Table~\ref{tab:s-targeting}), so most of what there is to reach
it reaches. The near-zero entry in the consistency table measures the yardstick,
not the score.

Against the exact gain the score does rank patients, for the modalities that
have per-patient signal to rank: Spearman \tgtSignalRhoMin--\tgtSignalRhoMax\
for the laboratory panel and the radiograph. For the electrocardiogram and
echocardiogram, whose true mean gain in this context is \tgtEcgTrueBits\ and
\tgtEchoTrueBits\ bits, it does not, and we report that rather than averaging
it into a headline: a score cannot order patients by a quantity that is
essentially constant at zero. In decision terms, a policy that buys a test for
the highest-scoring tenth of patients collects
\tgtSignalCaptureMin--\tgtSignalCaptureMax\% of the information an oracle
ordering would collect with the same budget, against
\tgtSignalChanceMin--\tgtSignalChanceMax\% for random ordering. The lift is
real and it is not large; it is the honest size of per-patient targeting at this
cohort size.

\subsection{Patient-level ordering reaches the same performance with fewer tests}

We simulated ordering policies honestly: a policy's prediction for a patient uses
only the modalities that policy bought, read off the cross-fitted probability for
exactly that subset. Comparators are chosen so the claim cannot be won cheaply --
ordering everything (the performance ceiling), ordering nothing beyond the free
baseline (the cost floor), each fixed single test, random ordering at matched
budget, and the standard clinical heuristic of testing the intermediate-risk band.

At matched discrimination (within 0.005 AUROC of testing everyone), the INFOGAIN
policy orders \redTestsAvoided\% fewer tests, saving \$\redCostSaved\ per patient
(\redCostSavedPct\%), at AUROC \redAuroc\ versus \redAurocAll\
(Fig.~\ref{fig:frontier}, Table~\ref{tab:policy}). The forgone information is
\redInfoGap\ bits, which by
Theorem~\ref{thm:omit} certifies a net-benefit loss of at most
\redCertifiedLoss\ at every threshold in the clinical window -- a guarantee that
holds independently of whether this particular sample happened to be favourable.
Decision curves for the reduced policy overlap those of full testing across the
threshold range (Fig.~\ref{fig:dca}).

The comparison against the risk-band heuristic is the informative one. It is a
genuinely good policy -- it captures most of the achievable saving -- and the
margin over it is the part of the result attributable to information theory
rather than to common sense.

\subsection{Sensitivity to how tests came to be ordered}

Estimating what a test \emph{would have} shown for a patient who did not receive
it is a counterfactual, and it is identified only under missingness at random
given the observed context, plus overlap. Neither is free.

We report the overlap diagnostic directly: the propensity to have received each
test is modelled from the baseline context, and the fraction of patients outside
the region of common support is reported per modality (Supplementary Table~\ref{tab:s-overlap}).
Patients in that region are flagged rather than scored, because for them the
counterfactual is an extrapolation, not an estimate.

For missingness \emph{not} at random we repeated the entire analysis on a cohort
in which ordering depends on the latent factors the ordering clinician can see
but the model cannot (Supplementary~\S\ref{sec:s-ident}). The direction of the bias is
informative: unmeasured indication inflates the apparent value of the tests that
were selectively ordered, and therefore makes the reduction claim conservative
rather than optimistic -- a policy trained under MNAR orders \emph{more} tests
than necessary, not fewer.

% ===================================================================== %
\section{Discussion}

The contribution is a change of question. ``Does adding this modality improve the
model?'' is a question about a cohort and a research paper. ``Given what I
already know about this patient, what will this test tell me, and is that worth
its cost?'' is a question about a patient and a decision, and it has a different
answer for each patient. We have shown that the second question is
well-posed, estimable, and translatable into the units of clinical decision
analysis by an identity rather than an analogy.

Three aspects seem to us to generalise beyond this dataset. First,
\emph{context-relative} decomposition sidesteps the partial-information-decomposition
impasse without giving up the vocabulary that made PID attractive: by fixing the
context to the tests actually in hand, redundancy and synergy become differences
of estimable quantities rather than solutions to an underdetermined axiom system.
Second, the restricted usable information $I^{[a,b]}$ is, we think, the quantity
that information-theoretic analyses of clinical data should report. Total mutual
information counts bits earned at thresholds no clinician acts on, and the
discrepancy is not a fixed overhead one could correct for: across our four
endpoints the redeemable share ranges from \restrictedShareMin\% to
\restrictedShareMax\%, and nothing in an unrestricted bit count says which end
of that range it sits at. Third, the safe-omission
certificate inverts the usual burden of proof: instead of demonstrating that a
reduced policy performed acceptably in one retrospective sample, it bounds what
the reduction could cost, a statement that survives the sample.

\paragraph{Limitations.} The first limitation is the data. Every number
reported here comes from the calibrated simulator of Methods~\ref{sec:sim}, not
from patients. That design lets us check the estimator against a decomposition
known in closed form -- which is how we established that overlap rather than
cohort size governs whether synergy is visible, a result the real data could not
have given us -- but it means the clinical findings below are properties of a
model we wrote. Which modality is redundant for which endpoint, and how much
testing can be avoided, are empirical claims that the MIMIC-IV extraction will
settle and this manuscript does not. What the simulator does establish is that
the machinery recovers the right answer when the right answer is known.

$\Vc$-information is relative to the predictor family
and the sample size; a stronger encoder or a larger cohort raises it, and
reporting the family is part of reporting the number. We give the learning-curve
extrapolation and the two-sided finite-sample bracket (Supplementary~\S\ref{sec:s-finite}) so that
readers can see how much of the gap is estimation and how much is genuine
absence of signal, but the estimate remains a lower bound. Our headline
modality representations are deliberately interpretable -- electrocardiogram
interval measurements and CheXpert labels rather than raw waveforms and pixels --
which makes the decomposition auditable by a clinician at the cost of some
usable information; the raw-signal path is implemented, and quantifying what
that costs requires the full waveform and image download (Supplementary~\S\ref{sec:s-raw}).
Relatedly, synergistic terms are the last to become estimable, and the binding
constraint is the number of patients who received both modalities rather than
the cohort size, so a study whose acquisition rates leave a small overlap will
report a synergistic test as independent. The error is conservative -- it
understates a test's value rather than overstating it, and so cannot manufacture
a reduction that is not there -- but it is an error, and Supplementary
Tables~\ref{tab:s-firstcorrect} and~\ref{tab:s-power} give the sample sizes involved. MIMIC-IV, the target cohort, is
a single academic centre in one country, its test-ordering culture is its own,
and the cost figures are US list prices; the frontier is therefore reported in full so that a reader with a
different fee schedule can read off their own operating point. Finally, a
retrospective simulation of a testing policy is not a trial. The certificate in
Theorem~\ref{thm:omit} bounds the information-theoretic cost of omission, not
the sociotechnical consequences of changing what clinicians are asked to do.

\paragraph{Outlook.} The natural next step is prospective: the per-patient score
is computable at the moment of the ordering decision, from data already in the
record, which makes a stepped-wedge evaluation feasible without any new
instrumentation. Sequential acquisition is implemented and is where redundancy
should pay off most: re-scoring after each result is what lets a policy see that
the second of two overlapping tests has stopped being worth anything \emph{for
this patient}, which a one-shot policy scored against the free baseline cannot.
Settings where the binding constraint is turnaround time or access rather than
money need only a different line in the cost model.

% ===================================================================== %
\section{Methods}

\subsection{Cohort, linkage and temporal gating}
\label{sec:cohort}

Encounters are drawn from MIMIC-IV v3.1 \citep{johnson2023mimiciv} for patients
aged $\ge$18. The prediction anchor (``index time'') is emergency-department
registration where available and hospital admission otherwise. Modalities are
linked by \texttt{subject\_id} across MIMIC-IV-ECG \citep{gow2023mimicecg},
MIMIC-CXR-JPG \citep{johnson2019mimiccxr}, MIMIC-IV-Note
\citep{johnson2023mimicnote} and, where present, MIMIC-IV-ED.

Temporal gating is enforced by an explicit policy object recorded in the cohort
metadata, and it has to distinguish two roles that a naive ``nothing after the
index time'' rule conflates. \emph{Context} modalities -- demographics, vital
signs, prior medications, prior notes -- must precede the index time; they are
what the clinician already holds. The \emph{orderable} modalities being scored,
however, are by definition acquired \emph{after} the decision to order them: an
electrocardiogram requested in the emergency department is performed minutes to
hours after registration. A strict zero look-ahead therefore excludes every
candidate test and silently turns the analysis into an analysis of nothing --
which is what it did until a dry-run test on synthetic MIMIC-shaped files caught
it. Orderable modalities may be acquired within a 6-hour window after the index,
and to keep that from leaking, every index-anchored outcome clock starts at the
\emph{end} of that window: ``ICU transfer within 48\,h'' means between 6 and 54
hours after registration, and a transfer during the acquisition window makes the
encounter non-evaluable rather than a positive. Look-back windows are 24\,h for
laboratory, electrocardiogram, radiograph and notes, and 168\,h for prior
prescriptions.

Discharge summaries are excluded for every endpoint anchored at admission; they
are written days later and are the single largest leakage source in multimodal
MIMIC work. Splits are grouped by patient, never by encounter, and the primary
analysis uses a temporal split on \texttt{anchor\_year\_group}, the closest
approximation to prospective validation available within one database.

\subsection{Endpoints}

\emph{30-day mortality}: death within 30 days of the index time, from the earlier
of in-hospital \texttt{deathtime} and the state-registry \texttt{dod}; encounters
whose observability window closes before the horizon are non-evaluable.
\emph{30-day heart-failure readmission}: unplanned readmission carrying a
heart-failure diagnosis within 30 days of discharge, following the CMS code
family; in-hospital deaths are non-evaluable. \emph{Acute kidney injury}: KDIGO
stage $\ge$1 within 7 days by the creatinine criteria \citep{kdigo2012aki} --
a rise of $\ge$0.3\,mg/dL in any rolling 48-hour window, or $\ge$1.5$\times$
baseline, baseline being the minimum in the preceding 7 days; chronic-dialysis
patients are non-evaluable. \emph{ICU transfer}: transfer to an intensive care
unit within 48\,h; patients already in an ICU at the index time are
non-evaluable. Non-evaluable encounters are coded missing, never zero.

\subsection{Modality representations}

Laboratory: a 35-analyte panel resolved against \texttt{d\_labitems} by label,
each contributing a winsorised value and an explicit measured/not-measured
indicator -- whether a test was ordered is itself informative and hiding it
inside an imputed value conflates two signals. Vital signs: emergency-department
triage observations. Medications: binary exposure flags for 18 therapeutic
classes in the prior week. Electrocardiogram: interval and axis measurements
derived from the MIMIC-IV-ECG machine measurements (heart rate, PR, QRS, QT,
QTc by Bazett, P/QRS/T axes), each with a measured indicator, plus 18 keyword
flags extracted from the cardiologist report. Chest radiograph: the 14 CheXpert
labels \citep{irvin2019chexpert} expanded over their four states (positive,
negative, uncertain, not mentioned -- averaging these would be meaningless) plus
view position. Notes: TF-IDF over 1--2 grams reduced by truncated SVD to 64
components, with vocabulary and projection fitted on training folds only.
Raw-waveform and raw-image encoders (1-D and 2-D residual networks) are
implemented for the comparison in Supplementary~\S\ref{sec:s-raw}.

\subsection{The predictive family}
\label{sec:family}

$\Vc$ is realised as one network that encodes each modality to a common width,
replaces the embedding of an absent modality by a learned per-modality
``absent'' token, and fuses the resulting tokens by a two-layer transformer
followed by an explicit second-order term $\sum_{m<m'} e_m \odot e_{m'}$, before
a multilayer-perceptron head; presence flags are also fed to the head, so
``not ordered'' and ``ordered and unremarkable'' are distinguishable. Training
samples masks from a distribution with explicit mass on the empty set (which
anchors every pointwise value), on the full set, and on the baseline, with
Bernoulli dropout elsewhere.

We chose attention fusion on the reasoning that synergistic information lives in
interactions between modalities, and that attention gives every modality a
direct path to every other where a concatenation head must find the interaction
inside a generic multilayer perceptron. Crossing the two switches on
\ablNCohorts\ independent cohorts shows the reasoning was wrong
(Supplementary Table~\ref{tab:s-ablation}). Attention and concatenation recover
the same share of the known synergy once the second-order term is present
(\ablAttnFmSyn\% against \ablConcatFmSyn\%, with between-cohort ranges that
overlap), and concatenation reaches it \ablSpeedup$\times$ faster. Attention
\emph{without} the second-order term is the worst arm of the four
(\ablAttnNoFmSyn\%, and it misreads the regime of both synergistic modalities
in both cohorts), while concatenation without it is unharmed
(\ablConcatNoFmSyn\%, every regime correct). The reason is a capacity bottleneck, not an
absence. Write the head-$h$ output for a query token attending over modality
tokens $1$ and $2$ as
\[
  o_h \;=\; a_h v_h(e_1) + (1-a_h) v_h(e_2)
       \;=\; v_h(e_2) + a_h\big(v_h(e_1)-v_h(e_2)\big),
  \qquad a_h=\sigma(z_h),
\]
with $z_h \propto q_h(e_1)^\top k_h(e_1) - q_h(e_1)^\top k_h(e_2)$. The value
path is linear in the value vectors, and $v_h(e_1)$, $v_h(e_2)$ each depend on
one modality alone. A cross-modal product therefore does exist -- $z_h$ is
bilinear in one modality's query and the other's key -- but the \emph{only}
route by which $e_1$ and $e_2$ jointly determine $o_h$ is the scalar $a_h$. Each
head supplies exactly one such channel, so the block's cross-modal interaction
capacity grows with the number of heads and not otherwise. Expanding
$a_h \approx a^0 + a^0(1-a^0)z_h$ makes the second constraint visible: the same
scalar multiplies $v_h(e_1)$ and $v_h(e_2)$, so within a head the cross-modal
terms cannot be given coefficients independent of the within-modality ones. The
factorization-machine term supplies the products directly and is therefore
repairing a bottleneck specific to attention, rather than adding a capability
concatenation lacks -- concatenation leaves the features adjacent for the head
to multiply, where the transformer pools them first.

This predicts that with no second-order term, recovery of a known synergistic
interaction should rise with head count and saturate. \texttt{head\_capacity\_%
sweep} tests it over $1,2,4,8,16$ heads on two cohorts under complete
acquisition, with a second arm holding the per-head dimension fixed, since at
width 48 the sixteen-head setting leaves three dimensions per head and a decline
would otherwise be ambiguous between saturation and starvation.

We report the transformer configuration because it produced every number here
and is at parity, not because it earns its cost; a replication should prefer
concatenation. Nothing above depends on the choice: $\IV$ is an infimum over
the family, so a weaker head only ever \emph{under}-reports information, and
architectural choices affect tightness and never validity.

\subsection{Estimation, cross-fitting and uncertainty}

All quantities derive from the pointwise value
$\pvi_i(S) = \log_2 f_S(y_i\mid x_i^S) - \log_2 f_\varnothing(y_i)$
\citep{ethayarajh2022pvi}, clipped to $\pm 8$ bits. Two properties are enforced
rather than assumed. Evaluation is always out of fold: $K$-fold cross-fitting
grouped by patient \citep{chernozhukov2018dml} means each patient is scored by
models that never saw them, and the union of folds covers the cohort, which is
what makes per-patient scores available for everyone rather than for a test
split. The null model is the training-fold base rate, the exact minimiser over
the constant sub-family, rather than the network's own all-absent member; the
gap between the two is reported because any slack in it inflates every level
estimate equally.

Differences are always estimated pairwise -- both terms share the patients and
the null, so differencing before averaging removes the patient-level variance
that dominates each term separately, tightening the standard error by roughly an
order of magnitude at realistic cohort sizes. Confidence intervals are BCa
bootstrap by default, with the finite-sample empirical-Bernstein interval
\citep{maurer2009empirical} reported for the headline claims; between-seed
variance is added to the sampling variance in both. Significance of a
conditional gain is additionally assessed by a paired sign-flip permutation
test. The learning bias is corrected by refitting at several training fractions
and extrapolating $\hat I(m) = I_\infty - am^{-\alpha}$ with $\alpha$ bounded to
$[0.15, 2]$; a power law rather than the classical $1/m$ series is used because
the dominant bias here is learning, not plug-in, and forcing $\alpha=1$
over-extrapolates interaction terms. Capacity is bounded from the trained
weights by the norm-based Rademacher bound of \citet{golowich2018size}, giving
the two-sided bracket of Theorem~S7.

\subsection{Per-patient expected information gain}
\label{sec:gain}

For patient $i$ with context $x_i^S$,
$\Delta_i(m\mid S) = \mathbb{E}_{x_m\sim q(\cdot\mid x_i^S)}
\KL\big(f_{S\cup m}(\cdot\mid x_i^S,x_m)\,\|\,f_S(\cdot\mid x_i^S)\big)$,
the expected information gain of Bayesian experimental design
\citep{lindley1956measure,chaloner1995bayesian} evaluated inside $\Vc$. It uses
neither the realised test result nor the outcome, so it is computable at the
bedside. The predictive distribution $q$ of the unacquired result is obtained by
a conditional bootstrap: results are resampled from patients who did receive the
test and who are nearest neighbours in the fused representation the outcome
model itself uses. This inherits the true joint structure of a high-dimensional
modality without a density model and degrades transparently -- if no similar
patient ever received the test, the overlap diagnostic reports it rather than
extrapolating. Donors are drawn only from the training folds of the model doing
the scoring. A Gaussian conditional sampler and a context-free marginal sampler
(a deliberately mis-specified control, isolating how much of the variation in
$\Delta_i$ comes from conditioning) are reported alongside.

\subsection{Decision-analytic translation and policy simulation}

Net benefit follows \citet{vickers2006decision}. Since Theorem~\ref{thm:identity}
requires nested posteriors, the coarse model is projected onto
$\mathbb{E}[\eta_2\mid \eta_1]$ by isotonic binning before the identity is
evaluated, and the calibration gap between model-implied and empirical net
benefit is reported with every decision curve
\citep{vancalster2019calibration}. Policies assign each patient a personalised
panel and are scored using the cross-fitted probability for exactly that subset.
The cost model keeps money, turnaround time and invasiveness as separate axes,
combined only at the point of use with weights the analysis sweeps over, because
which axis binds differs by setting.

\subsection{Simulation with exact ground truth}
\label{sec:sim}

$L$ binary latent factors $z$ drive Gaussian modality read-outs
$X_m = W_m z + \varepsilon_m$, with $W_m$ having zero columns for factors that
modality cannot see, and a logistic outcome depending on $z$ alone with main
effects and \emph{centred} interaction terms. Centring is what makes a synergy
term pure: the marginal effect of each factor on the logit is exactly zero, so
neither modality alone carries the signal. Because $z$ is finite and the
read-outs Gaussian, $P(Y=1\mid x_S)$ follows by enumeration for every subset and
every patient, so all population and per-patient quantities have a ground truth.
Test ordering depends on severity, missing at random given the free modalities by
default and not-at-random under the sensitivity setting. The generator emits the
same cohort object as the PhysioNet extraction, so the identical analysis code
runs on both.

\subsection{Reporting and reproducibility}

Reporting follows TRIPOD+AI \citep{collins2024tripodllm} where applicable. Every
number in this manuscript is a macro generated from the results tree by
\texttt{scripts/make\_paper\_numbers.py}, so the text cannot drift from the run
that produced it; unavailable quantities typeset as a visible marker rather than
a stale value. The full analysis is reproduced by
\texttt{python scripts/run\_all.py}. All theoretical claims are numerically
certified by \texttt{python -m infogain.theory.verify}
(all claims pass: \thyAllHold).

% ===================================================================== %
\section*{Data availability}

MIMIC-IV, MIMIC-IV-ECG, MIMIC-CXR-JPG and MIMIC-IV-Note are available from
PhysioNet to credentialed users who complete the required human-subjects
training and sign the data use agreement; see \texttt{docs/DATA\_ACCESS.md} for
the exact download commands and expected layout. No patient-level data are
redistributed here. The simulated cohort used for validation is regenerated
deterministically by \texttt{scripts/make\_synthetic\_cohort.py}.

\section*{Code availability}

All analysis code is in the accompanying repository under an MIT licence,
including the PhysioNet extraction, the estimators, the theorem checks and the
figure generation.

\bibliographystyle{unsrtnat}
\bibliography{refs}

% ===================================================================== %
\clearpage
\section*{Tables}

\begin{table}[htbp]\centering\small
\caption{\textbf{Context-relative decomposition across all four endpoints.}
For each modality: value alone on top of the free baseline ($I_m$), value once
every other test is in hand ($U$), and the split of their difference into
redundant and synergistic bits. $p$ is a paired sign-flip permutation test on
the conditional gain. The same pair of modalities changes regime between
endpoints, which is the paper's central empirical observation.}
\label{tab:decomposition}
\fittable{\begin{tabular}{llrrrrr}
\toprule
\textbf{endpoint} & \textbf{modality} & \textbf{marginal} & \textbf{conditional} & \textbf{redundant} & \textbf{synergistic} & \textbf{p} \\
\midrule
aki\_7d & labs & 0.1186 & 0.1118 & 0.0068 & 0.0000 & 0.0002 \\
aki\_7d & ecg & -0.0005 & -0.0014 & 0.0009 & 0.0000 & 1.0000 \\
aki\_7d & cxr & 0.0086 & 0.0015 & 0.0071 & 0.0000 & 0.0046 \\
aki\_7d & echo & -0.0005 & -0.0013 & 0.0009 & 0.0000 & 1.0000 \\
hf\_readmission\_30d & labs & 0.0510 & 0.0235 & 0.0274 & 0.0000 & 0.0002 \\
hf\_readmission\_30d & ecg & 0.0103 & 0.0145 & 0.0000 & 0.0042 & 0.0002 \\
hf\_readmission\_30d & cxr & 0.0313 & 0.0026 & 0.0287 & 0.0000 & 0.0002 \\
hf\_readmission\_30d & echo & 0.0087 & 0.0122 & 0.0000 & 0.0035 & 0.0002 \\
icu\_transfer\_48h & labs & 0.0135 & 0.0137 & 0.0000 & 0.0002 & 0.0002 \\
icu\_transfer\_48h & ecg & -0.0007 & -0.0003 & 0.0000 & 0.0003 & 0.8776 \\
icu\_transfer\_48h & cxr & 0.0182 & 0.0200 & 0.0000 & 0.0017 & 0.0002 \\
icu\_transfer\_48h & echo & -0.0002 & -0.0004 & 0.0003 & 0.0000 & 0.9838 \\
mortality\_30d & labs & 0.0262 & 0.0222 & 0.0039 & 0.0000 & 0.0002 \\
mortality\_30d & ecg & -0.0006 & 0.0133 & 0.0000 & 0.0139 & 0.0002 \\
mortality\_30d & cxr & 0.0065 & 0.0156 & 0.0000 & 0.0092 & 0.0002 \\
mortality\_30d & echo & -0.0002 & -0.0007 & 0.0004 & 0.0000 & 0.9966 \\
\bottomrule
\end{tabular}
}
\end{table}

\begin{table}[htbp]\centering\small
\caption{\textbf{Ordering policies at the primary endpoint.} Tests ordered per
patient, cost, discrimination, and the information given up relative to testing
everyone. The comparison that matters is against the risk-band heuristic at
matched budget, not against the extremes.}
\label{tab:policy}
\fittable{\begin{tabular}{lrrrrrr}
\toprule
\textbf{policy} & \textbf{tests per patient} & \textbf{mean cost} & \textbf{auroc} & \textbf{info gap vs all bits} & \textbf{nb@0.02} & \textbf{nb@0.05} \\
\midrule
order\_all & 4.0000 & 455.0000 & 0.7898 & 0.0000 & 0.0550 & 0.0382 \\
baseline\_only & 0.0000 & 0.0000 & 0.6630 & 0.0347 & 0.0524 & 0.0282 \\
fixed:cxr & 1.0000 & 45.0000 & 0.6938 & 0.0272 & 0.0526 & 0.0306 \\
fixed:ecg & 1.0000 & 25.0000 & 0.6625 & 0.0336 & 0.0523 & 0.0283 \\
fixed:echo & 1.0000 & 350.0000 & 0.6625 & 0.0343 & 0.0524 & 0.0281 \\
fixed:labs & 1.0000 & 35.0000 & 0.7578 & 0.0116 & 0.0539 & 0.0355 \\
risk\_band[0.40,0.90] & 2.0000 & 227.5000 & 0.7351 & 0.0150 & 0.0531 & 0.0336 \\
risk\_band[0.25,0.95] & 2.8000 & 318.5000 & 0.7613 & 0.0087 & 0.0537 & 0.0359 \\
risk\_band[0.50,0.95] & 1.8000 & 204.7500 & 0.7288 & 0.0147 & 0.0527 & 0.0320 \\
risk\_band[0.60,1.00] & 1.6000 & 182.0000 & 0.7228 & 0.0147 & 0.0525 & 0.0305 \\
random & 0.9987 & 112.9475 & 0.7005 & 0.0267 & 0.0529 & 0.0313 \\
random & 2.0070 & 228.2546 & 0.7313 & 0.0182 & 0.0535 & 0.0334 \\
random & 3.0047 & 341.8289 & 0.7602 & 0.0094 & 0.0542 & 0.0357 \\
\bottomrule
\end{tabular}
}
\end{table}

% ===================================================================== %
\clearpage
\section*{Figures}

\begin{figure}[htbp]\centering
\infogainfig[\textwidth]{fig1_decomposition}
\caption{\textbf{Context-relative decomposition of each modality's usable
information.} Blue: marginal value on top of the free baseline. Green: value
remaining once every other test is in hand. The hatched wedge is the difference,
blue where the context already contained those bits (redundant) and red where
they exist only in combination (synergistic). Bars are cross-fitted estimates
with 95\% intervals including between-seed variance. The accompanying CSV gives
the same numbers.}
\label{fig:decomp}
\end{figure}

\begin{figure}[htbp]\centering
\infogainfig[0.62\textwidth]{fig2_interaction_map}
\caption{\textbf{Modality-pair interaction map.} Signed interaction
$I(ab)-I(a)-I(b)$, measured on top of the free baseline. Blue is redundant (the
pair overlaps), red is synergistic (the pair is worth more than the sum of its
parts). A single signed number per pair requires no contested axiom to define.}
\label{fig:map}
\end{figure}

\begin{figure}[htbp]\centering
\infogainfig[\textwidth]{fig3_patient_gains}
\caption{\textbf{Who benefits from testing.} Left: distribution of the
per-patient prospective gain $\Delta_i(m\mid S_0)$ (log counts). Right: Lorenz
curves. Distance from the diagonal is the paper's central empirical claim -- if
benefit were uniform, patient-level targeting could not beat a blanket policy.}
\label{fig:gains}
\end{figure}

\begin{figure}[htbp]\centering
\infogainfig[0.62\textwidth]{fig4_decision_curves}
\caption{\textbf{Decision curves.} Net benefit against risk threshold for
testing everyone, testing no one beyond the free baseline, and the INFOGAIN
policy at its matched-performance operating point.}
\label{fig:dca}
\end{figure}

\begin{figure}[htbp]\centering
\infogainfig[0.66\textwidth]{fig5_reduction_frontier}
\caption{\textbf{Tests ordered against discrimination.} The INFOGAIN frontier
against every comparator: ordering everything, ordering nothing, each fixed
single test, random ordering at matched budget, and the risk-band heuristic. The
margin over the risk-band heuristic is the part attributable to information
theory rather than common sense.}
\label{fig:frontier}
\end{figure}

\begin{figure}[htbp]\centering
\infogainfig[0.62\textwidth]{fig7_identity}
\caption{\textbf{Theorem~\ref{thm:identity} made visible.} The information a
modality adds equals the area under its threshold-weighted net-benefit gain.
Shading marks the clinically usable window; bits outside it are earned at
thresholds nobody acts on.}
\label{fig:identity}
\end{figure}

\begin{figure}[htbp]\centering
\infogainfig[0.58\textwidth]{fig8_theorem1}
\caption{\textbf{Certified versus observed discrimination gain.}
Theorem~\ref{thm:auroc}'s lower bound on the achievable AUROC gain from adding
each modality, against the gain actually observed.}
\label{fig:thm1}
\end{figure}

\begin{figure}[htbp]\centering
\infogainfig[0.62\textwidth]{fig9_gain_by_risk}
\caption{\textbf{Expected gain against baseline risk.} Curves that were flat
would mean information-guided ordering reduces to risk stratification. They are
neither flat nor identically shaped across modalities.}
\label{fig:byrisk}
\end{figure}

\begin{figure}[htbp]\centering
\infogainfig[\textwidth]{figS1_recovery}
\caption{\textbf{Validation against known ground truth.} Left: estimated versus
true information at the largest simulated cohort, before and after
learning-curve correction. Right: fraction of the true information recovered as
a function of cohort size. The estimator approaches the truth from below, as a
computationally bounded quantity must.}
\label{fig:recovery}
\end{figure}

\end{document}
